\documentclass[11pt]{article}
\usepackage[margin=1in]{geometry}
\usepackage{amsmath,amssymb,amsthm,mathtools}
\usepackage{graphicx}
\usepackage{microtype}
\usepackage{xcolor}
\usepackage{hyperref}
\hypersetup{colorlinks=true,linkcolor=blue!55!black,citecolor=blue!55!black,urlcolor=blue!55!black}

\newtheorem{theorem}{Theorem}
\newtheorem{lemma}[theorem]{Lemma}
\newtheorem{corollary}[theorem]{Corollary}
\newcommand{\R}{\mathbb R}
\newcommand{\BMO}{\mathrm{BMO}}
\newcommand{\osc}{\mathcal O}
\newcommand{\E}{\mathbb E}
\newcommand{\1}{\mathbf 1}

\title{Radialization has BMO norm $\Theta(\sqrt n)$,\\
and symmetric rearrangement is not dimension-free}
\author{Solution packet for arXiv:2201.05130}
\date{17 August 2026}

\begin{document}
\maketitle

\begin{center}
\fbox{\parbox{0.91\textwidth}{\textbf{Status: substantial partial result, likely valid.}
The decreasing-rearrangement constant remains open.  The packet completely
classifies the dimension growth of the radialization step, proves that the
symmetric-rearrangement constant grows at least as $\sqrt n$, and improves its
known upper bound from exponential growth to $O(n)$.}}
\end{center}

\section{Source question and scope}

Burchard, Dafni, and Gibara ask for the sharp dimension dependence of the
decreasing-rearrangement constants $C_n$ and state that the sharp constants
and dimension dependence for symmetric decreasing rearrangement are also
unknown; see the introduction of arXiv:2201.05130, PDF pages 1--2
\cite{source}.  The full passage is reproduced below.

\begin{center}
\includegraphics[width=\textwidth]{figures/open_problem_crop.png}
\end{center}

For a cube $Q$ and a locally integrable $u$, write
\[
 \osc(u,Q)=\frac1{|Q|}\int_Q |u-u_Q|,\qquad
 \|u\|_{\BMO(\R^n)}=\sup_Q\osc(u,Q),
\]
where cubes have sides parallel to the coordinate axes.  Let $f^*$ be the
nonincreasing rearrangement of $|f|$, let $\omega_n$ be the volume of the
Euclidean unit ball, and set
\[
 Sf(x)=f^*(\omega_n|x|^n).
\]
Denote by $\mathfrak C_n$ and $\mathfrak D_n$ the optimal constants in
\[
 \|f^*\|_{\BMO(\R_+)}\le \mathfrak C_n\|f\|_{\BMO(\R^n)},
 \qquad
 \|Sf\|_{\BMO(\R^n)}\le \mathfrak D_n\|f\|_{\BMO(\R^n)}.
\]
The first source question---whether $\mathfrak C_n$ is dimension-free---is
not settled here.  The result answers the dimension-free part of the second
question negatively and isolates the precise growth of the intervening
radial operation.

\section{Proof intuition}

There are two independent mechanisms.  First, a decreasing one-dimensional
BMO function can change by at most a constant times its BMO norm whenever its
argument is multiplied by two.  Thus the oscillation of
$g(\omega_n|x|^n)$ on a cube is controlled by the oscillation of
$n\log|x|$.  For a uniform point $X$ in any cube, a Poincare estimate and a
dimension-sharp small-ball bound give
$\operatorname{Var}(\log|X|)=O(1/n)$.  Consequently radialization costs at
most $O(\sqrt n)$, and the logarithmic profile shows that this is sharp.

Second, to turn the logarithmic profile into the rearrangement of a function
whose original BMO norm stays bounded, use the cusp
\[
 f_n(x_1,x')=\bigl(-\log|x_1|-\|x'\|_\infty\bigr)_+.
\]
Its level sets have volume exactly $A_ne^{-t}$, so its symmetric
rearrangement is a multiple of $-n\log|x|$.  Yet on small cubes the two
terms in the cusp each have bounded oscillation, while on large cubes the
positive set is sufficiently thin to dilute the mean.  On a cube of the form
$a[1,2]^n$, the radial output is governed by
$\frac n2\log\sum_i X_i^2$; the central limit theorem produces a
$\sqrt n$ lower bound with an explicit constant.

\section{Main result}

For a nonnegative nonincreasing $g\in\BMO(\R_+)$ define
\[
 (T_ng)(x)=g(\omega_n|x|^n),
\]
and let $\mathfrak R_n$ be the operator norm of $T_n$ from this monotone cone
to $\BMO(\R^n)$.  Finally set
\[
 \kappa:=\frac1{7}\sqrt{\frac{17}{5\pi}}=0.148616\ldots.
\]

\begin{theorem}[Radialization and symmetric-rearrangement bounds]
For $n\ge3$,
\begin{equation}\label{eq:Rupper}
 \mathfrak R_n\le
 32\left(1+\frac{n\omega_n^{1/n}}{\pi\log2}
             \sqrt{\frac n{n-2}}\right)=O(\sqrt n),
\end{equation}
whereas
\begin{equation}\label{eq:Rlower}
 \liminf_{n\to\infty}\frac{\mathfrak R_n}{\sqrt n}
 \ge \frac{e\kappa}{4}>0.
\end{equation}
Hence $\mathfrak R_n=\Theta(\sqrt n)$.

Moreover, the optimal symmetric-rearrangement constants satisfy
\begin{equation}\label{eq:Dlower}
 \liminf_{n\to\infty}\frac{\mathfrak D_n}{\sqrt n}
 \ge \frac{\kappa}{4}>0,
\end{equation}
and
\begin{equation}\label{eq:Dupper}
 \mathfrak D_n\le 2(1+2\sqrt{n-1})\mathfrak R_n=O(n).
\end{equation}
In particular, symmetric decreasing rearrangement does not have a
dimension-free cube-BMO bound.  If $\sup_n\mathfrak C_n<\infty$, then
$\mathfrak D_n=\Theta(\sqrt n)$.
\end{theorem}

The factor $2(1+2\sqrt{n-1})$ in \eqref{eq:Dupper} is Theorem 1.1 of
arXiv:2011.09111 (PDF page 2) \cite{bounds}.  Its published corollary used a
comparison between cubes and annular sectors and gave the much larger upper
bound
$2(1+2\sqrt{n-1})n^{n/2}\omega_n$ \cite[PDF pp.~3--4]{bounds}.

\section{The radialization upper bound}

We first record a logarithmic modulus that uses only one-dimensional BMO.

\begin{lemma}[Dyadic modulus]\label{lem:dyadic}
Let $g$ be nonincreasing on $\R_+$ and let
$K=\|g\|_{\BMO(\R_+)}$.  At Lebesgue points $s,t>0$,
\[
 |g(s)-g(t)|\le 16K\left(1+\frac{|\log(s/t)|}{\log2}\right).
\]
\end{lemma}

\begin{proof}
It suffices first to compare $a$ and $2a$.  Let $U,V$ be independent and
uniform on $I=(a/2,5a/2)$.  On a set of probability $1/4$, $U$ lies in
$(a/2,a)$ and $g(U)\ge g(a)$; on another set of probability $1/4$, $V$
lies in $(2a,5a/2)$ and $g(V)\le g(2a)$.  Including the reversed event,
\[
 \E|g(U)-g(V)|\ge\frac18\bigl(g(a)-g(2a)\bigr).
\]
On the other hand,
$\E|g(U)-g(V)|\le2\osc(g,I)\le2K$.  Therefore
$g(a)-g(2a)\le16K$.  A dyadic telescoping from $s$ to $t$ proves the
claim.
\end{proof}

\begin{lemma}[Logarithmic radius in a cube]\label{lem:logradius}
Let $X$ be uniform in an arbitrary cube of side length $\ell$ in $\R^n$,
$n\ge3$.  Then
\begin{equation}\label{eq:logvar}
 \operatorname{Var}(\log|X|)
 \le \frac1{\pi^2}\frac n{n-2}\omega_n^{2/n}.
\end{equation}
\end{lemma}

\begin{proof}
The Neumann Poincare inequality on a cube gives, first for a smooth
regularization of the logarithm and then by approximation,
\[
 \operatorname{Var}(\log|X|)
 \le\frac{\ell^2}{\pi^2}\E|X|^{-2}.
\]
For every $r>0$, the elementary volume bound
\[
 \mathbb P\{|X|<r\}\le
 \min\left(1,\frac{\omega_n r^n}{\ell^n}\right)
\]
holds regardless of the location of the cube.  Put
$r_0=\ell\omega_n^{-1/n}$.  Layer-cake integration yields
\begin{align*}
 \E|X|^{-2}
 &=2\int_0^\infty r^{-3}\mathbb P\{|X|<r\}\,dr\\
 &\le \frac{2\omega_n}{\ell^n}\int_0^{r_0}r^{n-3}\,dr
       +2\int_{r_0}^\infty r^{-3}\,dr
 =\frac n{n-2}\frac{\omega_n^{2/n}}{\ell^2}.
\end{align*}
This proves \eqref{eq:logvar}.
\end{proof}

\begin{proof}[Proof of \eqref{eq:Rupper}]
Fix a cube $Q$ and let $X$ be uniform on it.  Set
$S=\omega_n|X|^n$ and $s_0=\exp(\E\log S)$.  Lemma
\ref{lem:dyadic}, followed by the elementary inequality
$\osc(Y)\le2\E|Y-c|$, gives
\begin{align*}
 \osc(T_ng,Q)
 &\le2\E|g(S)-g(s_0)|\\
 &\le32K\left(1+\frac1{\log2}\E|\log S-\E\log S|\right)\\
 &\le32K\left(1+\frac n{\log2}
        \sqrt{\operatorname{Var}(\log|X|)}\right).
\end{align*}
Insert Lemma \ref{lem:logradius} and take the supremum over $Q$.
Finally, Stirling's formula gives
$\omega_n^{1/n}\sim\sqrt{2\pi e/n}$, so the displayed bound is
$O(\sqrt n)$.
\end{proof}

\section{Sharpness of radialization}

We use the logarithmic profile
\[
 g_0(s)=(-\log s)_+.
\]
For completeness, $\|g_0\|_{\BMO(\R_+)}\le4/e$.  Indeed, truncation does
not increase mean oscillation.  If an interval $I$ of length $\ell$ lies at
distance at least $\ell$ from zero, then $-\log s$ varies by at most
$\log2$ on $I$.  Otherwise $I\subset(0,2\ell)$, the containing interval
has at most twice its length, and direct integration gives
$\osc(-\log s,(0,2\ell))=2/e$.

Now
\[
 T_ng_0(x)=\bigl(-\log(\omega_n|x|^n)\bigr)_+
          =n\log(R_n/|x|)_+,
 \qquad R_n=\omega_n^{-1/n}.
\]
Take $a_n=R_n/(3\sqrt n)$ and $Q_n=a_n[1,2]^n$.  This cube lies inside
$B(0,R_n)$.  If $X_1,\ldots,X_n$ are independent and uniform on $[1,2]$,
then addition of constants and scaling show
\begin{equation}\label{eq:Mn}
 \osc(T_ng_0,Q_n)
 =\E\left|\frac n2\log\Bigl(\sum_{i=1}^nX_i^2\Bigr)
 -\E\frac n2\log\Bigl(\sum_{i=1}^nX_i^2\Bigr)\right|.
\end{equation}
Let $Z_i=X_i^2$.  A direct calculation gives
\[
 \mu:=\E Z_i=\frac73,\qquad
 \sigma^2:=\operatorname{Var}(Z_i)=\frac{34}{45}.
\]
The delta method applied to $\sum_iZ_i$, together with boundedness of the
$Z_i$ (which supplies uniform integrability), implies that
\[
 \frac1{\sqrt n}\left[
 \frac n2\log\Bigl(\sum_i Z_i\Bigr)
 -\E\frac n2\log\Bigl(\sum_i Z_i\Bigr)\right]
 \longrightarrow N\left(0,\frac{\sigma^2}{4\mu^2}\right)
\]
in distribution and in first absolute moment.  Hence the expression in
\eqref{eq:Mn}, divided by $\sqrt n$, tends to
\[
 \frac{\sigma}{2\mu}\E|N(0,1)|
 =\frac{\sigma}{\mu\sqrt{2\pi}}
 =\kappa.
\]
Since $\|g_0\|_{\BMO}\le4/e$, this proves \eqref{eq:Rlower}.

\section{A uniformly BMO logarithmic cusp}

Define, for $n\ge2$ and $x'=(x_2,\ldots,x_n)$,
\begin{equation}\label{eq:cusp}
 f_n(x_1,x')=\bigl(-\log|x_1|-\|x'\|_\infty\bigr)_+.
\end{equation}

\begin{lemma}[Uniform source norm]\label{lem:cuspBMO}
$f_n$ is rearrangeable and
$\|f_n\|_{\BMO(\R^n)}\le4$ for every $n\ge2$.
\end{lemma}

\begin{proof}
We first use the dimension-free bound
\begin{equation}\label{eq:logBMO}
 \|\log|\,\cdot\,|\|_{\BMO(\R)}\le\frac8e.
\end{equation}
If an interval $I$ of length $\ell$ has distance at least $\ell$ from zero,
the oscillation is at most the range $\log2$.  Otherwise
$I\subset(-2\ell,2\ell)$.  Mean oscillation on a subset is at most the
volume ratio times the oscillation on the containing interval, while
\[
 \osc(-\log|t|,(-2\ell,2\ell))
 =\int_0^1|-\log s-1|\,ds=\frac2e.
\]
The volume ratio is $4$, proving \eqref{eq:logBMO}.

Let $Q=I_1\times\cdots\times I_n$ have side length $\ell$.  If
$\ell\le1$, truncation and the triangle inequality for mean oscillation
give
\[
 \osc(f_n,Q)
 \le \osc(-\log|x_1|,Q)+\osc(\|x'\|_\infty,Q)
 \le \frac8e+\ell<4,
\]
because $\|x'\|_\infty$ varies by at most $\ell$ on $Q$.

If $\ell>1$, nonnegativity gives
$\osc(f_n,Q)\le2|Q|^{-1}\int_Q f_n$.  For fixed $x'$ and
$v=\|x'\|_\infty$,
\[
 \int_{I_1}(-\log|x_1|-v)_+\,dx_1
 \le\int_{|x_1|<e^{-v}}(-\log|x_1|-v)\,dx_1
 =2e^{-v}.
\]
Consequently $|Q|^{-1}\int_Q f_n\le2/\ell$, and
$\osc(f_n,Q)<4$.  The level-set computation below also shows directly that
$f_n$ is rearrangeable.
\end{proof}

For every $t\ge0$, Fubini's theorem and the volume of an
$\ell_\infty$-ball give
\begin{align*}
 |\{f_n>t\}|
 &=2e^{-t}\int_{\R^{n-1}}e^{-\|x'\|_\infty}\,dx'\\
 &=2e^{-t}\int_0^\infty e^{-r}\,d(2r)^{n-1}
 =A_ne^{-t},
 \qquad A_n=2^n(n-1)!.
\end{align*}
It follows exactly that
\[
 f_n^*(s)=\log(A_n/s)_+,
 \qquad
 Sf_n(x)=n\log(\rho_n/|x|)_+,
 \qquad
 \rho_n=(A_n/\omega_n)^{1/n}.
\]
Taking the cube $\rho_n(3\sqrt n)^{-1}[1,2]^n$ repeats the calculation
\eqref{eq:Mn}; hence
\[
 \liminf_{n\to\infty}
 \frac{\|Sf_n\|_{\BMO}}{\sqrt n}\ge\kappa.
\]
Lemma \ref{lem:cuspBMO} proves \eqref{eq:Dlower}.

Finally, $Sf=T_n(f^*)$, so by definition
$\mathfrak D_n\le\mathfrak R_n\mathfrak C_n$.  Theorem 1.1 of
\cite{bounds} gives
$\mathfrak C_n\le2(1+2\sqrt{n-1})$, proving \eqref{eq:Dupper} and the
conditional $\Theta(\sqrt n)$ statement.

\section{Verification and limitations}

The analytic proof was checked in four independent pieces: the dyadic BMO
modulus, the layer-cake calculation in Lemma \ref{lem:logradius}, the exact
cusp distribution, and the delta-method constant.  The seeded script
\texttt{code/verify\_asymptotics.py} sampled $200{,}000$ vectors for each of
$n=10,25,50,100,200$.  The normalized values were respectively
\[
 0.15054,\ 0.14948,\ 0.14889,\ 0.14830,\ 0.14871,
\]
against the proved limit $\kappa=0.148616\ldots$.  This computation is a
sanity check only; the proof is the central limit theorem plus uniform
integrability.

A separate finite-grid search for a growing lower bound on
$\mathfrak C_n$ was rejected: ratios above one disappeared after arbitrary
shifted cubes were included.  It supplies no evidence or premise used here.
The present result leaves the sharp growth of $\mathfrak C_n$ open and leaves
a gap $\sqrt n\lesssim\mathfrak D_n\lesssim n$.  Closing the upper gap by
the factorization $Sf=T_n(f^*)$ is equivalent to obtaining a dimension-free
bound for $\mathfrak C_n$.

\paragraph{Bounded novelty search.}
On 17 August 2026, searches of arXiv/web used ids 2201.05130 and 2011.09111
and the phrases ``symmetric decreasing rearrangement BMO dimension lower
bound'', ``radially decreasing BMO cube norm'', ``logarithmic cusp'', and
``symmetrization BMO dimension-free''.  They returned the two cited source
papers and unrelated rearrangement literature, but no later match for
\eqref{eq:Rupper}--\eqref{eq:Dupper}.  This is bounded evidence, not an
exhaustive claim of novelty.

\paragraph{Human-review recommendation.}
Priority checks are: (i) the use of cube Poincare after regularizing
$\log|x|$ at the origin; (ii) the truncation and large-cube estimates in
Lemma \ref{lem:cuspBMO}; and (iii) uniform integrability in the delta-method
limit.  These are the only points where an omitted technical detail could
affect the theorem.  No computational dependency remains.

\begin{thebibliography}{9}
\bibitem{source}
A. Burchard, G. Dafni, and R. Gibara,
\emph{Vanishing Mean Oscillation and Continuity of Rearrangements},
arXiv:2201.05130 (2022), especially introduction, PDF pp.~1--2.

\bibitem{bounds}
A. Burchard, G. Dafni, and R. Gibara,
\emph{Mean Oscillation Bounds on Rearrangements},
Trans. Amer. Math. Soc. 375 (2022), 4429--4444;
arXiv:2011.09111, Theorems 1.1--1.2.
\end{thebibliography}

\end{document}
